\section{Evaluation}

\subsection{Demographic Distribution Analysis}

Both datasets exhibit a pronounced male skew across 200 professions, but the imbalance
is substantially more severe in SDv1.5 than in Re-LAION-5B. Re-LAION-5B contains
41.5\% female-labelled images against 58.5\% male-labelled, a gap of 17 percentage
points relative to parity. SDv1.5 widens this to 42 percentage points, with only
29.02\% female images. Comparing against BLS 2024 and ILOSTAT baselines, Re-LAION-5B
falls below real-world female proportions in the majority of ISCO groups relative to
US norms while remaining broadly consistent with global workforce proportions in
several others. SDv1.5 then amplifies this male skew further across nearly every
occupational group regardless of which baseline is used, confirming a two-stage bias
pattern in which the generative model exacerbates rather than merely inherits the
imbalances of its training corpus. Gender amplification is confirmed in 178 of 200
professions, with extreme cases including Logger (97.5\% male), Nurse (97.5\% female),
and Receptionist (97.3\% female) in SDv1.5.

The skin tone analysis reveals a near-complete collapse of dark-toned representation in
generative outputs. Both datasets are dominated by Mid-bin predictions, but the Dark-bin
proportion falls from 0.81\% in Re-LAION-5B to 0.14\% in SDv1.5, an 83\% reduction,
while the Light-bin proportion nearly doubles from 12.45\% to 22.87\%. Skin tone
amplification toward lighter tones is confirmed in 175 of 200 professions. A direct
comparison against real-world skin tone statistics was not conducted, as the
race-to-skin-tone reference mapping derived via FairFace collapsed all racial groups
predominantly to the Mid bin, providing insufficient discriminative resolution to serve
as a meaningful baseline.

\subsection{Dataset Classification}

A ConvNeXt-Tiny classifier trained on transformed image representations distinguishes
Re-LAION-5B, SDv1.5, and COCO with high accuracy across every transformation tested,
from unmodified images at 99.42\% down to mean RGB compression at 64.19\%. A
consistent hierarchy of discriminating power emerges: frequency and fine spatial
statistics (high-pass filter 98.64\%, patch shuffle $8{\times}8$ at 97.91\%) follow
closely behind the unmodified baseline; structural geometry transformations (depth
96.78\%, Canny edge 93.62\%) occupy the next tier; dense semantic representations
(segmentation 91.36\%, captions 90.99\%, object detection 88.62\%) follow; and global
colour statistics form the lowest tier (pixel shuffle 80.67\%, mean RGB 64.19\%).
Critically, no single transformation eliminates separability, confirming that dataset
identity is redundantly encoded across colour, texture, spatial structure, geometry,
semantics, and language simultaneously. VAE reconstruction reduces accuracy by only
0.27 percentage points relative to the unmodified baseline (99.15\% vs.\ 99.42\%),
establishing that low-level acquisition artefacts are not the dominant source of
discriminating signal. The substantially higher accuracies observed here relative to
the 82.0\% reference reported by Zeng et al.\ for purely web-scraped datasets are
consistent with SDv1.5's generative pipeline imposing consistent visual priors that
produce far greater internal homogeneity than a web-scraped corpus.

\subsection{Demographic Artifact Probing}

The occlusion experiments reveal a fundamental asymmetry between the two datasets.
In Re-LAION-5B, all person-hidden occlusion variants collapse to chance for both
gender and skin tone regardless of whether the background is retained or removed,
establishing that background context alone carries no recoverable demographic signal
in the web-scraped corpus. In SDv1.5, background context is a substantial independent
carrier of demographic signal: with the person masked and background retained, gender
AUC reaches 0.747 and skin tone AUC reaches 0.728. Silhouette information contributes
additively, with MaskSegm NoBg reaching gender AUC 0.834. This directly replicates the
contextual artifact identified by Meister et al.\ and demonstrates it is substantially
more pronounced in generated imagery than in web-scraped data.

Beyond occlusion, both datasets exhibit modest demographic signal in non-appearance
representations. Mean RGB achieves gender AUC 0.600/0.579 and skin tone AUC 0.563/0.543
for Re-LAION-5B and SDv1.5 respectively, consistent with the 0.580 reported by Meister
et al.\ on COCO. Pose keypoints achieve gender AUC 0.697/0.680 across datasets,
confirming that systematic body configuration differences contribute independently of
photometric properties. Object detection reveals gender-stereotyped scene compositions
at AUC 0.692 on Re-LAION-5B (0.608 restricted), while the gap narrows substantially
in SDv1.5. Caption probing reveals a cross-dataset reversal absent from all visual
probes: Re-LAION-5B achieves gender AUC 0.901 and skin tone AUC 0.749 in captions,
substantially exceeding SDv1.5 at 0.771 and 0.673 respectively. This inversion
indicates that web-scraped images carry richer linguistically recoverable demographic
associations, while SDv1.5 embeds stronger demographic signal in scene-level visual
statistics. Together, the two datasets exhibit qualitatively distinct spurious feature
profiles: Re-LAION-5B localises demographic signal within the person region, while
SDv1.5 distributes it pervasively across scene context, body geometry, object
composition, and linguistic register.

\subsection{Debiasing Evaluation}

The combined ITI-GEN, ControlNet, and Chamfer colour alignment pipeline produces a
measurable and substantive departure from the demographic and spurious feature profile
of standard SDv1.5 generation. The debiased corpus achieves the intended 50\% female /
50\% male distribution and a uniform spread across all ten MST levels, standing in
direct contrast to the standard pipeline's 29.02\% female and near-zero Dark-bin
representation.

Person-level probing confirms successful demographic conditioning: gender AUC reaches
$\approx$1.000 on unmodified images, as expected given that labels are derived from the
generation configuration. All non-person gender probes collapse to at or near chance:
mean RGB falls to 0.507, object detection to 0.449/0.457, depth at pose keypoints to
0.454, pose keypoints to 0.475, and caption probing to 0.530, a reduction of 24.1
percentage points from the standard SDv1.5 value of 0.771. For skin tone, all
non-person probes similarly collapse: mean RGB falls to 0.482, object detection to
0.439/0.461, depth at keypoints to 0.443, pose keypoints to 0.483, and captions to
0.490, all at or near the 10-class chance level of 0.500. The mean RGB skin tone result
of 0.482 provides the clearest mechanistic confirmation that the Chamfer colour
alignment successfully disrupts photometric spurious correlations at the global level.

Residual background-level signal persists in both analyses. The background-only gender
probe (MaskRect) achieves AUC 0.995 on the debiased corpus, and the skin tone
background-only probe retains 0.632, indicating that scene composition remains
partially entangled with demographic conditioning even under explicit pose and
appearance control. This represents the primary open challenge for future work on this
pipeline.